% Part: sets-functions-relations
% Chapter: sets
% Section: subsets

\documentclass[../../../include/open-logic-section]{subfiles}

\begin{document}

\olfileid[fa-IR]{sfr}{set}{sub}
\olsection{زیرمجموعه‌ها و مجموعه‌های توانی}

\begin{explain}
اغلب می‌خواهیم مجموعه‌ها را با یکدیگر مقایسه کنیم. یک مقایسهٔ آشکار
می‌تواند چنین باشد: \emph{هر آنچه در یک مجموعه است، در مجموعهٔ دیگر
نیز هست}. این وضعیت آن‌قدر مهم است که برای آن نمادگذاری تازه‌ای
معرفی کنیم.
\end{explain}

\begin{defn}[زیرمجموعه]
اگر هر !!{element} مجموعهٔ $A$، !!a{element} مجموعهٔ~$B$ نیز باشد،
می‌گوییم $A$ \emph{زیرمجموعهٔ}~$B$ است و می‌نویسیم $A \subseteq B$.
اگر $A$ زیرمجموعهٔ~$B$ نباشد، می‌نویسیم $A \not\subseteq B$.
اگر $A \subseteq B$ اما $A \neq B$، می‌نویسیم $A \subsetneq B$ و
می‌گوییم $A$ \emph{زیرمجموعهٔ سرهٔ} $B$ است.
\end{defn}

\begin{ex}
هر مجموعه زیرمجموعهٔ خودش است و $\emptyset$ زیرمجموعهٔ هر مجموعه است.
مجموعهٔ اعداد طبیعی زوج زیرمجموعهٔ مجموعهٔ اعداد طبیعی است. همچنین
$\{ a, b \} \subseteq \{ a, b, c \}$. اما $\{ a, b, e
\}$ زیرمجموعهٔ $\{ a, b, c \}$ نیست.
\end{ex}

\begin{ex}
عدد $2$ یک !!{element} مجموعهٔ اعداد صحیح است، درحالی‌که مجموعهٔ
اعداد زوج زیرمجموعهٔ مجموعهٔ اعداد صحیح است. بااین‌حال، ممکن است
مجموعه‌ای \emph{هم} !!a{element} و هم زیرمجموعهٔ مجموعه‌ای دیگر باشد؛
برای نمونه، $\{0\} \in \{0, \{0\}\}$ و نیز $\{0\} \subseteq \{0,
\{0\}\}$.
\end{ex}

اصل گسترش ملاکی برای یکسانی مجموعه‌ها به دست می‌دهد: $A = B$ اگر و
تنها اگر هر !!{element} مجموعهٔ~$A$، !!a{element} مجموعهٔ~$B$ نیز باشد
و برعکس. تعریف «زیرمجموعه» عبارت $A \subseteq B$ را دقیقاً به‌منزلهٔ
نیمهٔ نخست این ملاک تعریف می‌کند: هر !!{element} مجموعهٔ~$A$،
!!a{element} مجموعهٔ~$B$ نیز هست. البته اگر جای $A$ و $B$ را عوض کنیم،
این تعریف همچنان برقرار است؛ یعنی $B \subseteq A$ اگر و تنها اگر هر
!!{element} مجموعهٔ~$B$، !!a{element} مجموعهٔ~$A$ نیز باشد. این نیز
دقیقاً بخش «و برعکس» در اصل گسترش است. به بیان دیگر، از اصل گسترش
نتیجه می‌شود که دو مجموعه برابرند اگر و تنها اگر زیرمجموعهٔ یکدیگر باشند.

\begin{prop}
$A = B$ اگر و تنها اگر هم $A \subseteq B$ و هم $B \subseteq A$.
\end{prop}

اکنون فرصت مناسبی است تا چند نماد سودمند دیگر نیز معرفی کنیم. در تعریف
زیرمجموعه‌بودن $A$ از~$B$ گفتیم «هر !!{element} مجموعهٔ~$A$، \dots»
و جای «$\dots$» را با «!!a{element} مجموعهٔ $B$» پر کردیم. اما این
\emph{ساختار} عبارت چنان پرکاربرد است که بهتر است نمادگذاری صوری‌ای
برای آن معرفی کنیم.

\begin{defn}\ollabel{forallxina}
$(\forall x \in A)\phi$ صورت اختصاریِ $\forall x(x \in A \lif
\phi)$ است. به همین ترتیب، $(\exists x \in A)\phi$ صورت اختصاریِ
$\exists x(x \in A \land \phi)$ است.
\end{defn}

با این نمادگذاری می‌توانیم بگوییم $A \subseteq B$ اگر و تنها اگر
$(\forall x \in A)x \in B$.

اکنون به بررسی نوع خاصی از مجموعه می‌پردازیم: مجموعهٔ همهٔ
زیرمجموعه‌های یک مجموعهٔ داده‌شده.

\begin{defn}[مجموعهٔ توانی]
مجموعه‌ای که از همهٔ زیرمجموعه‌های مجموعهٔ~$A$ تشکیل شده است،
\emph{مجموعهٔ توانیِ}~$A$ نام دارد و آن را $\Pow{A}$ می‌نویسند.
  \[
    \Pow{A} = \Setabs{B}{B \subseteq A}
  \]
\end{defn}

\begin{ex}
همهٔ زیرمجموعه‌های ممکنِ $\{ a, b, c \}$ کدام‌اند؟ این زیرمجموعه‌ها
عبارت‌اند از: $\emptyset$، $\{a \}$، $\{b\}$، $\{c\}$، $\{a, b\}$،
$\{a, c\}$، $\{b,
c\}$، $\{a, b, c\}$. مجموعهٔ همهٔ این زیرمجموعه‌ها
$\Pow{\{a,b,c\}}$ است:
\[
\Pow{\{ a, b, c \}} = \{\emptyset, \{a \}, \{b\}, \{c\}, \{a, b\},
\{b, c\}, \{a, c\}, \{a, b, c\}\}
\]
\end{ex}

\begin{prob}
همهٔ زیرمجموعه‌های $\{a, b, c, d\}$ را فهرست کنید.
\end{prob}

\begin{prob}
نشان دهید اگر $A$ دارای $n$ !!{element}s باشد، آنگاه $\Pow{A}$ دارای
$2^n$ !!{element}s است.
\end{prob}

\end{document}
